\section{The Key You Already Designed}
\label{sec:ch08-the-key-you-already-designed}
\index{global object identification!as a federation key|(}%

\code{@key(fields: "id")} names a field, and in Mosaic that field is not a
\code{Guid}. Chapter~\ref{ch:schema-design-that-survives-change} made
\code{Product.id} a global object identifier, so it answers with an opaque
string carrying the type name beside the key:

\begin{minted}[fontsize=\footnotesize]{json}
{"data":{"products":[
  {"id":"UHJvZHVjdDoAAACgAAAAQIAAAAAAAAAB","sku":"MOS-FRN-0001","title":"Larsen Oak Dining Table"},
  {"id":"UHJvZHVjdDoAAACgAAAAQIAAAAAAAAAC","sku":"MOS-FRN-0002","title":"Larsen Dining Chair"}]}}
\end{minted}

Base64, decoding to \code{Product:} followed by the sixteen bytes of
\code{a0000000-0000-4000-8000-000000000001}. That string is now the federation
key, which means it is what travels in every representation, and it means the
two services have to agree about it exactly.

\index{global object identification!type name inside the key}%
They do agree, and the reason is the part of the encoding that looks like
decoration. The GraphQL type name is inside the key. Both services call the type
\code{Product}, both use HotChocolate's default node identifier serialiser, so
both turn the same \code{Guid} into the same string without ever having spoken.
Neither of those two conditions is enforced anywhere: the serialiser is
configurable, and a key is a string that either side can produce by any means it
likes. I have not built the graph where they disagree, so I will not tell you
what it looks like from the outside. I will say that a composer compares
schemas, and two schemas can be identical while the strings behind them are not.
Figure~\ref{fig:ch08-key-round-trip} is the round trip.

\begin{figure}[htbp]
  \centering
  % What happens to a product identifier between the two services: encoded on the
% way out by [ID], decoded on the way in by hand. Referenced from chapter 08.
% Bare tikzpicture; the figure environment, caption and label live at the call
% site.
\begin{tikzpicture}[
    font=\small,
    box/.style={draw, rounded corners=2pt, minimum width=24mm,
                minimum height=11mm, align=center, font=\scriptsize},
    wide/.style={draw, rounded corners=2pt, minimum width=58mm,
                 minimum height=11mm, align=center, font=\scriptsize},
    flow/.style={-{Stealth[length=2mm]}},
    lbl/.style={font=\scriptsize, align=center, inner sep=1pt},
    note/.style={font=\scriptsize, align=center, gray, inner sep=1pt}
  ]

  \node[box]  (out) at (1.2,0)  {Catalog holds\\\code{a0000000-...-0001}};
  \node[wide] (enc) at (6.35,0) {\code{UHJvZHVjdDoAAACgAAAAQIAAAAAAAAAB}};
  \node[box]  (in)  at (11.5,0) {Mosaic wants\\\code{a0000000-...-0001}};

  \draw[flow] (out) -- (enc);
  \draw[flow] (enc) -- (in);

  \node[lbl,  anchor=south] at (2.93,0.75) {\code{[ID]} encodes};
  \node[note, anchor=south] at (2.93,1.20) {automatic};

  \node[lbl,  anchor=south] at (9.77,0.75) {\code{Parse(id, typeof(Guid))}};
  \node[note, anchor=south] at (9.77,1.20) {yours to write};

  \node[note, anchor=north, align=center] at (6.35,-0.75)
    {the same sixteen bytes throughout, with the GraphQL type name in front of
    them:\\that prefix is the only thing distinguishing this from a customer's
    identifier};

\end{tikzpicture}

  \caption{One identifier, three forms. Encoding happens on the way out of
  either service and HotChocolate does it; decoding happens on the way in and
  nobody does it for you. The string in the middle is what a representation
  carries, and the type name inside it is what makes two processes that have
  never met agree on a key.}
  \label{fig:ch08-key-round-trip}
\end{figure}

\subsection*{Nothing decodes it for you}

\index{federation!and relay identifiers}%
The federation package and the global object identification feature do not know
about each other at all. Grep the whole federation package's source at tag
\code{16.6.0} for \code{NodeId} or for \code{Relay} and you get nothing back.

A reference resolver's key parameter is read out of the representation by the
federation argument parser, which coerces the literal with the field's own
scalar and then converts the result to the parameter's type.
The scalar is \code{ID}, so the coercion hands back the base64 string unchanged,
and the conversion to \code{Guid} fails. When it fails, the parser returns the
default. So a reference resolver written the obvious way:

\begin{minted}{csharp}
[ReferenceResolver]
public static Product? ResolveReference(Guid id) => ...
\end{minted}

is handed \code{Guid.Empty} for a perfectly valid key, looks it up, finds
nothing, and answers:

\begin{minted}{json}
{"data":{"_entities":[null]}}
\end{minted}

\index{federation!silent null entity}%
Not an error. A null, in the slot where the router expected a product, for every
representation. And the specification's \code{[\_Entity]} return type is
nullable so that a subgraph can answer null for one entity without failing the
batch~\autocite{apollo2026subgraphspec}, which means nothing anywhere in the
system considers this a fault.

\code{[ID]} on the parameter does not help. I tried it, because it is the first
thing anybody tries. An entity whose reference resolver was declared
\code{ResolveByKey([ID] Guid id)} resolved correctly when handed a raw
\code{Guid} string and answered null when handed the encoded key, which is the
exact opposite of what the attribute promises. Argument coercion is where
\code{[ID]} does its work, and a representation never goes through argument
coercion.

\subsection*{The decode, and two things about it that cost time}

\index{INodeIdSerializer@\code{INodeIdSerializer}}%
So the decode is explicit, and it lives in one small class per service:

\begin{minted}[fontsize=\footnotesize]{csharp}
public static bool TryDecode(string formattedId, IResolverContext context, out Guid id)
{
    var services = context.Schema.Services;
    var accessor = services?.GetService(typeof(INodeIdSerializerAccessor))
        as INodeIdSerializerAccessor;

    if (accessor is null)
    {
        id = Guid.Empty;
        return false;
    }

    try
    {
        var nodeId = accessor.Serializer.Parse(formattedId, typeof(Guid));

        if (nodeId.TypeName != nameof(Product) || nodeId.InternalId is not Guid decoded)
        {
            id = Guid.Empty;
            return false;
        }

        id = decoded;
        return true;
    }
    catch (Exception)
    {
        id = Guid.Empty;
        return false;
    }
}
\end{minted}

\medskip\noindent\textbf{The second argument is the type of the identifier, not
of the entity.} Pass \code{typeof(Guid)} and it works. Pass
\code{typeof(Product)} and it throws, with a message reporting that no
serialiser is registered for a type you can see in your own schema, which is a
confusing five minutes. This is not a guess about the intended usage: it is the
call HotChocolate itself makes for an ordinary \code{[ID]} argument, in
\code{GlobalIdInputValueFormatter}, where the runtime type passed in is the type
of the identifier value.

\medskip\noindent\textbf{The serialiser is a schema service, not a request
service.} Ask \code{context.Services} for it and you get null; ask
\code{context.Schema.Services} and you get it. That is why the method takes an
\code{IResolverContext} rather than the accessor itself, and it is the same
constraint section~\ref{sec:ch08-two-attributes} described from the other side:
a schema service cannot be a reference resolver parameter, so it has to be
reached through one that can.

\medskip
The type name check earns its place, and not in the way I expected when I wrote
it. Hand Catalog a customer key with and without that line and the answer is the
same both times: a null. Mosaic's identifiers are \code{Guid}s, so a customer's
decodes to a value no product has ever had, and the lookup misses either way.

It matters the moment the internal identifier is not globally unique, which is
the ordinary case rather than the exotic one. A graph keyed on integers encodes
customer 5 and product 5 into two different strings that decode to the same
\code{5}, and a resolver that ignores the type name would answer with product 5
for a representation naming a customer. The check is the difference between a
service that knows which type it was asked about and one that trusts a number.
Mosaic gets away without it; I would not want the next service to inherit the
habit.

\subsection*{Two copies of one file, on purpose}

\index{extraction!contracts between services}%
\code{ProductKey} exists twice, once in each service, with the same method body
and a different comment above it. I considered a shared package and decided
against it, and I would decide the same way again: the moment the two services
share a class, the format of the key becomes a compile-time dependency between
them, and the fact that they have to \emph{agree about a string} disappears
into a project reference. Two services that must agree about a format are
better off saying so twice than importing it once.

That is also the honest version of what the extraction actually did to
\code{Product}. There are two \code{Product} types in this repository now, in
two assemblies, with different fields. They are the same GraphQL type because
they carry the same name and the same key, and for no other reason at all.

\index{global object identification!as a federation key|)}%
